\documentclass[11pt,a4paper]{article}

\usepackage[T1]{fontenc}
\usepackage[utf8]{inputenc}
\usepackage[ngerman]{babel}

\usepackage{amsmath,amssymb,amsfonts,amsthm}
\usepackage{mathtools}
\usepackage{graphicx}
\usepackage{geometry}
\usepackage{hyperref}
\usepackage{enumitem}
\usepackage{physics}
\usepackage{bm}

\geometry{margin=2.5cm}

\title{
	Holonomie, EABC-Strukturen und ein spektrales Modell\\
	für Teilchen, Familien und emergente Wechselwirkungen\\[0.5em]
	\large Ein Forschungsprogramm zwischen Algebra,
	Topologie, Zahlentheorie und mathematischer Physik
}

\author{
	Thomas Hoffbauer
}

\date{\today}

\newtheorem{definition}{Definition}
\newtheorem{satz}{Satz}
\newtheorem{lemma}{Lemma}
\newtheorem{korollar}{Korollar}
\newtheorem{bemerkung}{Bemerkung}
\newtheorem{beispiel}{Beispiel}

\begin{document}
	
	\maketitle
	
	\begin{abstract}
		
		Die vorliegende Arbeit entwickelt ein spektrales Forschungsprogramm,
		das diskrete Zahlentheorie, Gruppenstrukturen,
		Topologie und geometrische Modelle miteinander verbindet.
		
		Ausgangspunkt ist die Beobachtung,
		dass die Einheitengruppe modulo $12$
		
		\[
		(\mathbb Z/12\mathbb Z)^\times
		=
		\{1,5,7,11\}
		\]
		
		eine natürliche Viererstruktur besitzt,
		die sowohl algebraisch als auch geometrisch
		durch einen regulären Tetraeder dargestellt werden kann.
		
		Auf dieser Basis wird ein Modell entwickelt,
		in welchem Spin, Familienstruktur,
		Holonomie, Neutrinooszillation,
		Van-der-Waals-Wechselwirkungen
		sowie spektrale Stabilitätsphänomene
		als unterschiedliche Manifestationen
		einer zugrundeliegenden Holonomiestruktur erscheinen.
		
		Historische Bezüge reichen von
		Pythagoras und Platon
		über Hamilton, Klein, Poincaré,
		Dirac, Aharonov und Bohm
		bis hin zu von Klitzing, Onsager
		und modernen Konzepten der
		nichtkommutativen Geometrie.
		
		Das Modell erhebt keinen Anspruch,
		eine abgeschlossene physikalische Theorie darzustellen,
		sondern formuliert ein mathematisches Forschungsprogramm,
		dessen zentrale Objekte
		diskrete Holonomien,
		Windungszahlen
		und spektrale Operatoren sind.
		
	\end{abstract}
	
	\tableofcontents
	
	\newpage
	
	%%%%%%%%%%%%%%%%%%%%%%%%%%%%%%%%%%%%%%%%%%%%%%%%%%%%%%%%%%
	\section{Historische Motivation}
	%%%%%%%%%%%%%%%%%%%%%%%%%%%%%%%%%%%%%%%%%%%%%%%%%%%%%%%%%%
	
	Die Geschichte der Physik zeigt,
	dass fundamentale Fortschritte häufig dann entstanden,
	wenn scheinbar voneinander unabhängige mathematische Strukturen
	als verschiedene Erscheinungsformen
	einer tieferen Ordnung erkannt wurden.
	
	Die Geometrisierung der Mechanik durch Newton,
	die Vereinheitlichung von Elektrizität und Magnetismus
	durch Maxwell,
	die Raumzeitgeometrie Einsteins,
	die Spinorentheorie Diracs
	sowie die topologische Quantisierung
	im Quanten-Hall-Effekt
	gehören zu den bekanntesten Beispielen.
	
	Parallel dazu entwickelte sich in der Mathematik
	eine zweite Tradition.
	
	Von den pythagoreischen Zahlenfiguren
	über Platons reguläre Körper,
	die Quaternionen Hamiltons,
	die Gruppentheorie Kleins,
	die Topologie Poincarés,
	die Spektraltheorie Hilberts
	bis zur nichtkommutativen Geometrie Connes'
	taucht immer wieder dieselbe Frage auf:
	
	\begin{quote}
		Welche Strukturen bleiben erhalten,
		wenn lokale Details verschwinden?
	\end{quote}
	
	Die moderne Topologie beantwortet diese Frage
	durch Homotopie, Homologie und Holonomie.
	
	Die vorliegende Arbeit verfolgt die Hypothese,
	dass genau diese Sichtweise
	auch für diskrete arithmetische Strukturen fruchtbar sein könnte.
	
	%%%%%%%%%%%%%%%%%%%%%%%%%%%%%%%%%%%%%%%%%%%%%%%%%%%%%%%%%%
	\section{Die EABC-Struktur}
	%%%%%%%%%%%%%%%%%%%%%%%%%%%%%%%%%%%%%%%%%%%%%%%%%%%%%%%%%%
	
	Betrachtet wird die Einheitengruppe modulo $12$
	
	\[
	(\mathbb Z/12\mathbb Z)^\times
	=
	\{1,5,7,11\}.
	\]
	
	Wir schreiben
	
	\[
	E\equiv1,
	\qquad
	A\equiv5,
	\qquad
	B\equiv7,
	\qquad
	C\equiv11.
	\]
	
	Es gilt
	
	\[
	A^2=B^2=C^2=E.
	\]
	
	Ferner
	
	\[
	AB=C,
	\qquad
	AC=B,
	\qquad
	BC=A.
	\]
	
	Die Struktur ist isomorph zur Kleinschen Vierergruppe
	
	\[
	V_4
	\cong
	(\mathbb Z/2\mathbb Z)^2.
	\]
	
	Historisch erscheint dieselbe Struktur
	bereits implizit bei Klein,
	explizit in der modernen Gruppentheorie
	sowie in vielen Klassifikationen diskreter Symmetrien.
	
	%%%%%%%%%%%%%%%%%%%%%%%%%%%%%%%%%%%%%%%%%%%%%%%%%%%%%%%%%%
	\section{Geometrische Realisierung}
	%%%%%%%%%%%%%%%%%%%%%%%%%%%%%%%%%%%%%%%%%%%%%%%%%%%%%%%%%%
	
	Wir definieren
	
	\[
	E=(1,1,1),
	\]
	
	\[
	A=(-1,-1,1),
	\]
	
	\[
	B=(1,-1,-1),
	\]
	
	\[
	C=(-1,1,-1).
	\]
	
	Dann folgt unmittelbar
	
	\[
	E+A+B+C=0.
	\]
	
	Die vier Punkte bilden einen regulären Tetraeder.
	
	Die quadratische Kantenlänge beträgt
	
	\[
	s^2=8.
	\]
	
	Diese Zahl wird später
	als fundamentaler Energiesprung interpretiert.
	
	%%%%%%%%%%%%%%%%%%%%%%%%%%%%%%%%%%%%%%%%%%%%%%%%%%%%%%%%%%
	\section{Quaternionische Hebung}
	%%%%%%%%%%%%%%%%%%%%%%%%%%%%%%%%%%%%%%%%%%%%%%%%%%%%%%%%%%
	
	Die EABC-Struktur besitzt eine natürliche Hebung
	zur Quaternionenalgebra
	
	\[
	\mathbb H.
	\]
	
	Wir identifizieren
	
	\[
	E \mapsto \pm1,
	\]
	
	\[
	A \mapsto \pm i,
	\]
	
	\[
	B \mapsto \pm j,
	\]
	
	\[
	C \mapsto \pm k.
	\]
	
	Wesentlich ist
	
	\[
	V_4
	\simeq
	Q_8/\{\pm1\}.
	\]
	
	Damit entsteht erstmals Orientierung.
	
	Insbesondere gilt
	
	\[
	BC=+A,
	\]
	
	aber
	
	\[
	CB=-A.
	\]
	
	Die Reihenfolge wird physikalisch relevant.
	
	Hier erscheint zum ersten Mal
	eine mathematische Analogie zum Spinbegriff Diracs.
	
	%%%%%%%%%%%%%%%%%%%%%%%%%%%%%%%%%%%%%%%%%%%%%%%%%%%%%%%%%%
	\section{Holonomie als zentrales Prinzip}
	%%%%%%%%%%%%%%%%%%%%%%%%%%%%%%%%%%%%%%%%%%%%%%%%%%%%%%%%%%
	
	Der entscheidende Gedanke dieser Arbeit lautet:
	
	\begin{quote}
		Nicht Zustände,
		sondern Umläufe sind fundamental.
	\end{quote}
	
	Für einen geschlossenen Pfad $P$
	definieren wir
	
	\[
	\Gamma(P)
	=
	\prod_{(ij)\in P}
	U_{ij}.
	\]
	
	Die Holonomie ist die diskrete Entsprechung
	der kontinuierlichen Phase
	
	\[
	\exp\left(
	i\oint A_\mu dx^\mu
	\right).
	\]
	
	Damit entsteht eine direkte Analogie
	zum Aharonov--Bohm-Effekt.
	
	%%%%%%%%%%%%%%%%%%%%%%%%%%%%%%%%%%%%%%%%%%%%%%%%%%%%%%%%%%
	\section{Aharonov--Bohm, von Klitzing und Onsager}
	%%%%%%%%%%%%%%%%%%%%%%%%%%%%%%%%%%%%%%%%%%%%%%%%%%%%%%%%%%
	
	Diese drei Nobelpreis-gekrönten Entdeckungen
	bilden den konzeptionellen Kern des Modells.
	
	\subsection{Aharonov--Bohm}
	
	Die beobachtbare Größe ist nicht das lokale Feld,
	sondern die Umlaufphase.
	
	\[
	\Gamma
	=
	e^{i\phi}.
	\]
	
	\subsection{Von Klitzing}
	
	Im Quanten-Hall-Effekt erscheint eine
	topologische Quantenzahl
	
	\[
	\nu.
	\]
	
	Im vorliegenden Modell wird
	
	\[
	\nu(P)
	=
	\sum m_{ij}
	\pmod{24}
	\]
	
	als Windungszahl interpretiert.
	
	\subsection{Onsager}
	
	Für den inversen Pfad gilt
	
	\[
	\Gamma(P^{-1})
	=
	\Gamma(P)^{-1}.
	\]
	
	Die reziproke Größe lautet
	
	\[
	\Omega(P)
	=
	\frac12
	\left(
	\Gamma(P)
	+
	\Gamma(P^{-1})
	\right).
	\]
	
	also
	
	\[
	\Omega(P)
	=
	\cos\phi.
	\]
	
	Damit erscheinen Holonomie,
	Topologie und Reziprozität
	als verschiedene Aspekte desselben Objekts.
	
	%%%%%%%%%%%%%%%%%%%%%%%%%%%%%%%%%%%%%%%%%%%%%%%%%%%%%%%%%%
	\section{Spektrale Dynamik}
	%%%%%%%%%%%%%%%%%%%%%%%%%%%%%%%%%%%%%%%%%%%%%%%%%%%%%%%%%%
	
	Der fundamentale Operator lautet
	
	\[
	H_{BM}
	=
	\Delta_{\mathcal G}
	+
	\alpha \log G
	+
	\beta \|S\|^2
	+
	\gamma \Omega(K)
	+
	\lambda \|L\|^2
	+
	\chi(\Gamma).
	\]
	
	Die Eigenwertgleichung
	
	\[
	H_{BM}\psi
	=
	E\psi
	\]
	
	liefert die zulässigen Zustände.
	
	Damit wird das gesamte Modell
	von einer statischen Struktur
	zu einer spektralen Theorie.
	
	%%%%%%%%%%%%%%%%%%%%%%%%%%%%%%%%%%%%%%%%%%%%%%%%%%%%%%%%%%
	\section{Diracoperator}
	%%%%%%%%%%%%%%%%%%%%%%%%%%%%%%%%%%%%%%%%%%%%%%%%%%%%%%%%%%
	
	Analog zu Dirac wird ein Operator eingeführt
	
	\[
	D_{BM}
	=
	iT_A
	+
	jT_B
	+
	kT_C
	+
	\omega T_\Gamma.
	\]
	
	mit
	
	\[
	\omega
	=
	e^{2\pi i/24}.
	\]
	
	Die zentrale Hypothese lautet
	
	\[
	H_{BM}
	=
	D_{BM}^2.
	\]
	
	Spin erscheint damit
	als Konsequenz der Quaternionenstruktur.
	
	%%%%%%%%%%%%%%%%%%%%%%%%%%%%%%%%%%%%%%%%%%%%%%%%%%%%%%%%%%
	\section{Ausblick}
	%%%%%%%%%%%%%%%%%%%%%%%%%%%%%%%%%%%%%%%%%%%%%%%%%%%%%%%%%%
	
	Das hier vorgestellte Programm
	ist keine fertige physikalische Theorie.
	
	Vielmehr formuliert es eine Reihe mathematischer Hypothesen:
	
	\begin{enumerate}
		\item Holonomien sind fundamentaler als Teilchen.
		\item Familien entsprechen Windungszuständen.
		\item Neutrinooszillationen sind Rotationen im Holonomieraum.
		\item Van-der-Waals-Wechselwirkungen entstehen durch
		Synchronisation von Holonomieschwankungen.
		\item Teilchenmassen entsprechen Eigenwerten
		eines diskreten Spektraloperators.
	\end{enumerate}
	
	Die zentrale Forschungsfrage lautet:
	
	\[
	\boxed{
		\text{Ist die Physik das Spektrum eines diskreten Holonomieoperators?}
	}
	\]
	
\end{document}